The central inquiry guiding this dissertation is how acts of modification within digital audio practice can generate new forms of knowledge grounded in culture and embodiment. From this overarching concern emerge the following questions:

\subsection{Primary Research Question}

\noindent How can sustained engagement with digital audio tools, guided by Black epistemologies, lead to new approaches in sonic interaction design and tool development?

\subsection{ Supporting Questions}

\noindent In what ways do acts of modification function as forms of \textit{transduction}, converting speculative cultural imagination into technical and material form?

\medskip

\noindent How can individual practitioners systematically identify and document the cultural assumptions embedded in digital creation tools' interface design, algorithmic processes, and interaction paradigms?

\medskip

\noindent How can procedural systems, interactive media platforms, and computational creative tools be modified to encode Black spatial, temporal, and embodied epistemologies?

\medskip

\noindent What systematic documentation and annotation practices enable practitioners to preserve and transmit cultural decision-making processes within tools?

\medskip

\noindent How do these modifications alter perception and creative practice? What kinds of \textit{refraction} occur when culturally informed tools reshape the act of listening, making, and collaboration?

\medskip

\noindent How might a framework provide a methodological bridge between critical theory and practice, transforming critique into processes of collaborative and embodied making?